\chapter{自我 $\Psi$ 的引入}
\chaptervoice{自我不是贴在记忆外面的姓名，也不是永不变化的人格文件；它是贯穿当前状态、经历选择与生成结构的共同约束，使变化仍能属于同一条可追溯轨迹。}

\section*{理论来路与依赖接口}
第25章已经给出带结构时间的生命周期，但 $h$、$E$、$\Theta$ 本身只说明状态、轨迹与生成结构，尚未决定哪些资源属于本主体、哪些承诺必须延续、哪些经验与长期价值相关。《比较智能差异.md》早期曾把自我放在 $\Theta$ 之后作为更慢的第四层，后段又明确修正：$\Psi$ 不是第四种记忆，而是同一自我结构在三相中的投影。本章采用这一较成熟定义，同时保留自我的慢时间尺度。

令
\begin{equation}
 \Psi_t=(I_t,V_t,C_t,O_t,R_t),\qquad
 \Psi_t\xrightarrow{\Pi_h,\Pi_E,\Pi_\Theta}
 (\Psi_h,\Psi_E,\Psi_\Theta),
\end{equation}
其中 $I$ 是身份谱系，$V$ 是价值与禁限，$C$ 是持续承诺，$O$ 是所有权和权限边界，$R$ 是关系与角色模型。三种投影不是三个自我，而是同一约束在不同数学对象中的作用接口。

\section{为什么 $h/E/\Theta$ 仍然缺少主体}
一套系统可以有活跃上下文、完整日志和强生成模型，却仍不能回答“谁对这次行动负责”。原因是三相记忆只给出信息的时间组织，不自动给出归属、承诺和授权。两套副本若共享同一 $\Theta$，但拥有不同密钥、债务、用户授权和后续承诺，它们不能仅因模型相同而被视作同一主体。

最低主体判据可写为约束集合
\begin{equation}
 \mathcal C(\Psi_t)=\{\text{lineage},\text{commitment},\text{authority},
 \text{value-boundary},\text{accountability}\}.
\end{equation}
它必须实际改变召回、行动或更新的可行集；只生成“我是某某”的文本而不改变任何决策，不构成功能性自我。反过来，系统即使没有自然语言自述，只要能持续维护授权、承诺、谱系和责任，也可以具有本章意义上的主体中心。这里讨论的是工程上的持续主体，不由此推出主观体验。

\section{$\Psi$ 作为全息中心}
“全息”不是说每个记忆片段都完整包含全部自我，而是说同一身份约束在三个相态中可被局部读取并保持相容。对任意允许状态，投影需满足软一致性
\begin{equation}
 d_h(\widehat\Psi_h,\Pi_h\Psi)+d_E(\widehat\Psi_E,\Pi_E\Psi)
 +d_\Theta(\widehat\Psi_\Theta,\Pi_\Theta\Psi)\le \epsilon,
\end{equation}
但不要求三者数据格式相同。例如当前状态中的权限令牌、经历中的签名主体和结构中的长期角色先验都指向同一谱系，却承担不同功能。

若三个投影冲突，不能以多数表决自动修复：缓存中的旧角色可能过期，日志中的历史授权也不代表当前仍有效。系统应按权威源、版本和有效期解析，并把冲突升级给治理层。因而“中心”是一组跨相态一致性关系，不是数据库中心表，也不是神经解剖位置；“全息”只是对分布式投影的结构比喻。

\section{$\Psi$ 在 $h$ 中的背景引力}
$\Psi_h$ 通过先验、注意偏置和行动约束调制当前状态。可将未调制激活记为 $\widetilde h_t$，则
\begin{equation}
 h_t=\operatorname{Project}_{\mathcal H(\Psi_t,b_t)}
 \bigl(\widetilde h_t+A_t\Pi_h(\Psi_t)\bigr),
\end{equation}
其中投影保证权限与安全边界，$A_t$ 决定当前任务下哪些身份成分相关。“背景引力”因此指持久而非独占的偏置：它改变信息进入工作集和候选行动的概率，却不能压制新的现实证据。

可通过保持观察与任务不变、受控替换 $\Psi_h$ 的角色或承诺成分，检验其是否改变注意、风险选择与资源分配。若自我先验总能否决反证，系统会形成自证循环；所以现实残差和硬安全信号必须有独立入口。适当的自我约束使当下行为连续，过强则变成僵化，过弱则使每次上下文都像一个无历史的新主体。

\section{$\Psi$ 在 $E$ 中的经历筛选}
$\Psi_E$ 决定哪些事件值得进入候选写入、以何种优先级召回，而不决定事实是否发生。写入门可分解为
\begin{equation}
 q_i=w_nN_i+w_gG_i+w_cC_i+w_rR_i,
\end{equation}
其中新颖性 $N_i$、目标相关性 $G_i$、承诺相关性 $C_i$ 与风险 $R_i$ 共同决定保留等级。安全事故即使不符合自我叙事，也必须由风险通道强制留存；原始审计记录不能被价值筛选删除。

召回同样需要双通道：与当前自我相关的利用检索，以及不受自我偏好支配的反证检索。否则系统会只记住支持既有身份的经历，逐步制造叙事闭环。评价 $\Psi_E$ 时，应测量关键承诺事件的召回率、反例覆盖率、来源完整性和存储成本，而不是只看个性化程度。筛选的对象是注意与维护预算，不是真实本身。

\section{$\Psi$ 在 $\Theta$ 中的结构空洞}
“结构空洞”表示许多知识、技能和关系以一个未被单条内容填满的主体索引为参照：谁能行动、谁承担结果、哪些边界不可交换。数学上可把 $\Theta$ 的部分关系写成带主体槽位的函数
\begin{equation}
 f_k(x;\Theta,\Psi)=f_k(x;\Theta,\operatorname{role}_k(\Psi)),
\end{equation}
其中槽位可由当前主体状态实例化，但不等于在 $\Theta$ 中复制一份完整 persona。删除该索引后，知识可能仍在，却无法稳定绑定所有权、第一人称承诺和权限。

“空洞”也不是缺失数据或神秘本体。它可由消融验证：若移除主体索引后，跨任务的责任绑定和承诺履行显著恶化，而一般事实预测近似不变，则其中心功能可辨识。若所有能力完全不受影响，所谓结构中心可能只是命名冗余。实现上要避免单点故障：谱系与权限可以分布式承载，但必须存在一致性和恢复协议。

\section{Who Am I}
“我是谁”不是请求模型即兴写自传，而是对当前可验证自我状态的查询。响应应由证据支持：谱系来自签名和检查点，能力来自近期评测，承诺来自权威账本，关系来自有时效的交互记录。可以定义自我报告
\begin{equation}
 S_t=\operatorname{Summarize}(\Psi_t;\mathcal Q),
 \qquad S_t=(\text{claim},\text{source},\text{validity},\text{confidence})_{1:m}.
\end{equation}
查询范围 $\mathcal Q$ 决定报告粒度；系统必须能说“不知道”或“需要重新验证”。

语言一致不代表身份连续，身份连续也不保证自我报告准确。一个固定提示词可以长期说同一句话，却没有共同谱系；同一持续系统也可能因状态估计错误而误报能力。因此需要分别测试叙述一致性、事实校准和控制效应。只有自我描述能被来源核验，并对后续授权与行动产生合规影响时，它才是功能自表示，而非风格文本。

\section{Self 作为跨尺度慢约束}
$\Psi$ 的典型时间尺度应慢于 $h$ 和普通 $\Theta$ 更新，但“慢”不是永不改变。连续近似可写为
\begin{equation}
 \tau_\Psi\dot\Psi
 =\Pi_{T_{\mathcal C_\Psi}(\Psi)}Q(E,\Theta,\Psi),
 \qquad \tau_h\ll\tau_E\ll\tau_\Theta\ll\tau_\Psi,
\end{equation}
其中切空间投影只允许满足身份、安全与承诺约束的局部漂移。工程实现通常是混杂过程：证据积累连续，价值或权限变更经审批后离散提交，并保留版本、迁移和回滚路径。

更新太快会使单次诱导改变长期身份；完全冻结又会阻断纠错、成长和角色退出。合理门槛依变更影响面而不同：能力估计可较快校准，长期承诺与根本价值必须需要更长证据窗口和更高权限。时间尺度链是设计假设，需由扰动实验测定，不能把某个固定比率当作自然定律。

\section{身份连续性}
身份稳定不是 $\Psi_t=\Psi_{t+1}$，而是相邻版本间存在保持关键不变量的有向迁移。令 $m_{t\to t'}$ 为迁移证明，则
\begin{equation}
 \operatorname{Cont}(t,t')=1
 \iff \exists m_{t\to t'}:
 \begin{cases}
 I_{t'}\sim I_t,\\
 C_{t'}=\operatorname{Resolve}(C_t,\Delta C),\\
 O_{t'}=\operatorname{Authorize}(O_t,\Delta O),\\
 d_V(V_t,V_{t'})\le\delta_V
 \end{cases}
\end{equation}
并且每项变化都有来源、批准者和现实后果。连续性因此是语义—因果谱系，不依赖组成硬件或模型参数逐位相同。

复制会揭示边界：两个实例可共享复制前历史，但复制后形成分叉谱系，不能同时把同一排他权限当作自己的。合并则必须显式解决冲突承诺，不能靠拼接记忆宣称成为一个主体。崩溃恢复若从受信检查点和事件日志重建，可保持工程连续；若来源不可验证，只能视为候选继承者。以上是 AgentOS 的责任与控制原则，不直接裁决哲学人格同一性。

\section*{本章收束：自我是约束关系，不是第四个仓库}
$\Psi$ 贯穿 $h$、$E$、$\Theta$：在 $h$ 中调制当前注意和可行动域，在 $E$ 中分配写入与召回预算，在 $\Theta$ 中提供主体索引与长期边界。它通过慢、可审计的漂移维持身份连续，却必须接受现实反证。下一章的 $\Omega$ 将描述生命周期的长期生成方向；它不能与回答“我是谁”的 $\Psi$ 混同。

\begin{readercheck}
为什么共享同一模型参数的两个副本不必是同一主体？$\Psi_E$ 可以调整什么、绝不能改写什么？请给出一次身份更新所需的谱系、承诺、权限与价值连续性证据。
\end{readercheck}
\cleardoublepage
